\documentclass[14pt]{extarticle}
\usepackage{fontspec}
\usepackage{polyglossia}
\usepackage{amsmath}
\usepackage{amsfonts}
\usepackage{amssymb}
\usepackage{pgfplots}
\usepackage{romannum}
\pgfplotsset{compat=1.18}

% Hebrew setup
\setdefaultlanguage{hebrew}
\setotherlanguage{english}
\newfontfamily\hebrewfont[Script=Hebrew]{Arial}

% Math font setup for proper Greek symbols
\usepackage{unicode-math}
\setmathfont{Latin Modern Math}

% Increase line spacing throughout document
\linespread{1.3}

\usepackage[margin=1in]{geometry}

% Left-aligned equation environment
\newenvironment{lefteqs}{%
    \begin{english}%
    \begin{flushleft}%
    \renewcommand{\arraystretch}{1.3}%
    $\begin{array}{@{}l@{}}%
}{%
    \end{array}$%
    \end{flushleft}%
    \end{english}%
}

\begin{document}

\begin{center}
    {\LARGE \textbf{מטלה 3 - אלקטרומגנטיות אנליטית}}\\[0.5em]
    {\large שי רואימי - 318986270}\\[0.3em]
\end{center}

\hrule
\vspace{2em}

\textbf{\large שאלה 1}

(פרק 2 שאלה 3)

בריק נתון $\vec{E_1}=-15\hat{x} + 20\hat{y}+30\hat{z}$, וגם $\epsilon_2=4\epsilon_0$.

\textbf{א)} 
נמצא את $\vec{D_2} \, , \vec{E_2}$, בהנחה שהמישור נטול מטען חופשי.

תנאי השפה הם רציפות $E$ הטנגנטיאלי, וכיוון ואין מטענים חופשיים על השפה אז $\vec{D_{1,z}=D_{2,z}}$ (רציפות $D$ הנורמלי).

\begin{lefteqs}
    \vec{E_{1,z}} = \vec{E_{2,z}} \Rightarrow \vec{E_2} = -15\hat{x} + 20\hat{y} + E_{2,z}\hat{z} \\
    \vec{D_{1,z}} = \vec{D_{2,z}} \Rightarrow \vec{D_z} = \epsilon \vec{E_z} \Rightarrow \vec{D_{1,z}} = \epsilon_0 E_{1,z} = \epsilon_2 E_{2,z} \\
    \epsilon_0 \cdot 30 = 4\epsilon_0 E_{2,z} \Rightarrow E_{2,z} = \frac{30}{4} = 7.5 \\
    \vec{E_2} = -15\hat{x} + 20\hat{y} + 7.5\hat{z} \\
    \vec{D_2} = \epsilon_2 \vec{E_2} = 4\epsilon_0(-15\hat{x} + 20\hat{y} + 7.5\hat{z}) \\
    \vec{D_2} = \epsilon_0(-60\hat{x} + 80\hat{y} + 30\hat{z})
\end{lefteqs}

\textbf{ב)} 
נמצא את הזווית $\theta$ שהשדה החשמלי $E_2$ יוצר עם הניצב למישור $z=0$.

\begin{lefteqs}
    \tan(\theta) = \frac{\sqrt{(-15)^2 + 20^2}}{7.5} \Rightarrow \theta = 73.3^\circ
\end{lefteqs}

\textbf{ג)} 
כעת נניח צפיפות מטען חופשית של $2 \cdot 10^{-10} \, \frac{C}{m^2}$ על המישור.

הרכיבים המקבילים של $E_2$ נותרו זהים, כלומר $E_{2,x} = E_{1,x}$ ו-$E_{2,y} = E_{1,y}$.

כיוון ויש מטען חופשי על המישור, אז $D_{1,z} - D_{2,z} = \sigma = 2 \cdot 10^{-10}$.

\begin{lefteqs}
    D_{1,z} = \epsilon_0 E_{1,z} \Rightarrow D_{2,z} = D_{1,z} - \sigma = \epsilon_0 \cdot 30 - 2 \cdot 10^{-10} \\
    \Rightarrow D_{2,z} = 6.55 \cdot 10^{-11} \\
    E_{2,z} = \frac{D_2,z}{\epsilon_2} = \frac{6.55 \cdot 10^{-11}}{4 \cdot 8.85 \cdot 10^{-12}} = 1.85 \\
    \vec{E_2} = -15\hat{x} + 20\hat{y} + 1.85\hat{z} \\
    \vec{D_2} = \epsilon_2 \vec{E_2} = 4\epsilon_0(-15\hat{x} + 20\hat{y} + 1.85\hat{z}) \\
    \vec{D_2} = \epsilon_0(-60\hat{x} + 80\hat{y} + 7.4\hat{z})
\end{lefteqs}

\vspace{3em}
\hrule
\vspace{2em}

\textbf{\large שאלה 2}

(פרק 2 שאלה 5)
מצא את השדה החשמלי בכל מקום עבור שני האלקטרטים הבאים:
\textbf{א)} 

כדור בעל רדיוס $a$ עם קיטוב $\vec{P} = P_0 \cdot \frac{r}{a} \hat{r}$

\begin{lefteqs}
    \rho_{pol} = -\vec{\nabla} \cdot \vec{P} = \frac{-1}{r^2} \frac{\partial}{\partial r} (r^2 P_r) = \frac{-1}{r^2} \frac{\partial}{\partial r} (r^2 P_0 \cdot \frac{r}{a}) \\
    \rho_{pol} = -\frac{P_0}{a} \cdot \frac{1}{r^2} \cdot 3r^2 = -\frac{3P_0}{a} \\
    \sigma_{pol} |^{r=a} = \vec{P} \cdot \hat{n} = \vec{P} \cdot \hat{r} = P_0 \cdot \frac{a}{a} = P_0 \\
    \underline{r<a:} \\
    Q_{enc} = \int_V \rho_{pol} dV = \int_0^r -\frac{3P_0}{a} \cdot 4\pi r\prime^2 \cdot dr\prime \\
    Q_{enc} = -\frac{12 \cdot P_0 \cdot \pi}{a} \left[ \frac{r\prime^3}{3} \right]_0^r = -\frac{4P_0\pi r^3}{a} \\
\end{lefteqs}

ע"פ חוק גאוס:

\begin{lefteqs}
    E \cdot 4\pi r^2 = \frac{Q_{enc}}{\epsilon_0} = \frac{-\frac{4P_0\pi r^3}{a}}{\epsilon_0} \\
    \vec{E} = -\frac{P_0 r}{a \epsilon_0} \hat{r} \\
    \underline{r>a:} \\
    Q_{enc}\prime = Q_{volume}|_{r=a} + Q_{surface}|_{r=a} = -\frac{4P_0\pi a^3}{a} + 4\pi a^2 P_0 \\
    \Rightarrow Q_{enc}\prime = 4\pi P_0 a^2 - 4\pi P_0 a^2 = 0 \Rightarrow \vec{E} = 0 \\
    E = \begin{cases}
    -\frac{P_0 r}{a \epsilon_0} \hat{r} & r < a \\
    0 & r > a
    \end{cases}
\end{lefteqs}

\textbf{ב)} 

לוח בעל עובי $b$, בעל קיטוב $P_0 \hat{x}$ הנמצא בין שני לוחות קבל שהם בהפרש פוטנציאל $V_0$.
נניח כי הקיטוב אני משתנה בין השדה החשמלי.

\begin{lefteqs}
    \rho_{pol} = -\vec{\nabla} \cdot \vec{P} = -\frac{\partial P_x}{\partial x} = -\frac{\partial}{\partial x} (P_0) = 0 \\
    \sigma_{pol} = \vec{P} \cdot \hat{n} \\
    \sigma_{top} = P_0 \hat{x} \cdot \hat{x} = P_0 \\
    \sigma_{bottom} = P_0 \hat{x} \cdot (-\hat{x}) = -P_0 \\
\end{lefteqs}

בציור נראה כי המרחק בין לוחות הקבל הוא $d$ אז ראשית:

\begin{lefteqs}
    \forall x>d, x<0 \, , E=0
\end{lefteqs}

נסתכל על $E$ הנוצר בתוך הלוח בעל עובי $b$, למעלה לו ומתחתיו.

הפנים התחתונות של הלוח בעלות צפיפות מטען $-P_0$, כלומר השדה מעל הפנים התחתונות הוא $-\frac{P_0}{2\epsilon_0}$.

הפנים העליונות, למעלה $\frac{P_0}{2\epsilon_0}$ ומתחת $\frac{P_0}{2\epsilon_0}$.

לכן בין הלוח לבין לוחות הקבל $E=0$ (כתוצאה מהלוח המקוטב), ובתוך האלקטרט השדה הוא $E_{bound} = \frac{-P_0}{\epsilon_0}$.

נסמן את השדה הנוצר מלוחות הקבל בתור $E_c$:

\begin{lefteqs}
    \int_0^d E_{total} dl = V_0 \Rightarrow \int_0^d (E_c + E_{bound}) dl
\end{lefteqs}

נניח שהאלקטרט במרכז בין לוחות הקבל:

\begin{lefteqs}
    = \int_{\frac{d+b}{2}}^d E_c dl + \int_{\frac{d-b}{2}}^{\frac{d+b}{2}} (E_c + E_{bound}) dl + \int_0^{\frac{d-b}{2}} E_c dl \\
    = E_c(d - \frac{d}{2} - \frac{b}{2}) + E_c(\frac{d}{2} - \frac{b}{2}) + (E_c + E_{bound})(\frac{d+b}{2} - \frac{d-b}{2}) \\
    = E_c(d-b) + (E_c - \frac{P_0}{\epsilon_0})\cdot b = V_0
    \Rightarrow E_c \cdot d - E_c \cdot b + E_c \cdot b -\frac{P_0}{\epsilon_0} \cdot b = V_0 \\ \\
    \Rightarrow E_c = \frac{V_0 + \frac{P_0 b}{\epsilon_0}}{d} = \frac{V_0}{d} + \frac{P_0 b}{\epsilon_0 d}
\end{lefteqs}

לסיכון, מעבר ללוחות הקבל $E=0$

בין לוחות הקבל לאקלטרט: $\vec{E} = (\frac{V_0}{d} + \frac{P_0 b}{\epsilon_0 d}) \hat{x}$

ובתוך האלקטרט: $\vec{E} = (\frac{V_0}{d} + \frac{P_0 b}{\epsilon_0 d} - \frac{P_0}{\epsilon_0}) \hat{x}$


\vspace{3em}
\hrule
\vspace{2em}

\textbf{\large שאלה 3}

\textbf{א)}
משוואה 2.14 בעמוד 138 היא:

\begin{lefteqs}
    Q_{pol} = \int_V \rho_{pol} dV + \int_S \sigma_{pol} dS = 0 \\
    \sigma_{pol} = \vec{P} \cdot \hat{n} \\
    \rho_{pol} = -\vec{\nabla} \cdot \vec{P}
\end{lefteqs}

ע"פ משפט הדיבגרנץ: $\iiint_V (\vec{nabla} \cdot \vec{F}) dV = \iint_S \vec{F} \cdot \hat{n} dS$.

\begin{lefteqs}
    \iiint_V \rho_{pol} dV = \iiint_V -\vec{\nabla} \cdot \vec{P} dV = -\iint_S \vec{P} \cdot \hat{n} dS = -\iint_S \sigma_{pol} dS \\
    Q_{pol} = -\iint_S \vec{P} \cdot \hat{n} dS + \iint_S \vec{P} \cdot \hat{n} dS = 0
\end{lefteqs}

\textbf{ב)}
המטען המושרה על השפה הוא:
\begin{lefteqs}
    \sigma_b = P_{R,z}|_{z=0} - P_{L,z}|_{z=0} = 0 - 0 = 0 \\
    \vec{P} = (\epsilon - \epsilon_0)\vec{E} \Rightarrow P_{R,z} = (\epsilon_1 - \epsilon_0)E_z^R \\
    P_{L,z} = (\epsilon_2 - \epsilon_0)E_z^L \\
    E_z^R|_{z=0} = -\frac{1}{4\pi\epsilon_1} \cdot \frac{d(Q' - Q)}{(r^2 + d^2)^{\frac{3}{2}}} \\
    E_z^L|_{z=0} = \frac{1}{4\pi\epsilon_2} \cdot \frac{d(Q'' d)}{(r^2 + d^2)^{\frac{3}{2}}} \\
    \sigma_b = -\frac{\epsilon_1 - \epsilon_0}{4\pi\epsilon_1} \cdot \frac{d(Q' - Q)}{(r^2+d^2)^{\frac{3}{2}}} - \frac{\epsilon_2 - \epsilon_0}{4\pi\epsilon_2} \cdot \frac{dQ''}{(r^2 + d^2)^{\frac{3}{2}}} \\
    Q' = \frac{\epsilon_1 - \epsilon_2}{\epsilon_1 + \epsilon_2} Q \\
    Q'' = \frac{2\epsilon_2}{\epsilon_1 + \epsilon_2} Q \\
    \sigma_b = -\frac{\epsilon_1 - \epsilon_0}{4\pi\epsilon_1} \cdot \frac{d(\frac{\epsilon_1 - \epsilon_2}{\epsilon_1 + \epsilon_2} Q - Q)}{(r^2 + d^2)^{\frac{3}{2}}} - \frac{\epsilon_2 - \epsilon_0}{4\pi\epsilon_2} \cdot \frac{d(\frac{2\epsilon_2}{\epsilon_1 + \epsilon_2} Q)}{(r^2 + d^2)^{\frac{3}{2}}} \\
\end{lefteqs}

\vspace{3em}
\hrule
\vspace{2em}

\textbf{\large שאלה 4}

נתון וקטור קיטוב של קוביה בראשית הצירים בעלת אורך $a$:

\begin{lefteqs}
    \vec{P}(\vec{r}) = P_0 \hat{x} + (\frac{z^2}{a^4} - \frac{x}{2a})P_0 \cdot \hat{z}
\end{lefteqs}

\textbf{א)}

נחשב את צפיפות המטען המשטחי המושרה על כל אחת מדפנות הקוביה ואת צפיפות המטען הנפחי המושרה בתוך הקוביה.

\underline{עבור הפאה $z=0$:}
\begin{lefteqs}
    \sigma_b = P(z=0) \cdot (-\hat{z}) = (P_0 \hat{x} - \frac{x}{2a}P_0\hat{z})(-\hat{z}) \\
    \sigma_b = \frac{x}{2a}P_0
\end{lefteqs}

\underline{עבור הפאה $x=0$:}
\begin{lefteqs}
    \sigma_b = P(x=0) \cdot (-\hat{x}) = (P_0 \hat{x} + \frac{z^2}{a^4}P_0\hat{z})(-\hat{x}) \\
    \sigma_b = -P_0
\end{lefteqs}

\underline{עבור הפאה $x=a$:}
\begin{lefteqs}
    \sigma_b = P(x=a) \cdot \hat{x} = (P_0 \hat{x} + (\frac{z^2}{a^4} - \frac{1}{2})P_0\hat{z})\hat{x} \\
    \sigma_b = P_0
\end{lefteqs}

\underline{עבור הפאה $z=a$:}
\begin{lefteqs}
    \sigma_b = P(z=a) \cdot \hat{z} = (1 - \frac{x}{2a})P_0 \\
\end{lefteqs}

\underline{עבור הפאות $y=0$ ו-$y=a$:}
\begin{lefteqs}
    \sigma_b = P(y=0) \cdot \hat{y} = P(y=a) \cdot (-\hat{y}) = 0
    -\vec{\nabla} \cdot \vec{P} = -\frac{\partial P_x}{\partial x} - \frac{\partial P_z}{\partial z} \\
    = -\frac{P_z}{\partial z} = -\frac{\partial}{\partial z} \left[ (\frac{z^4}{a^4} - \frac{x}{2a})P_0 \right] = -\frac{4z^3P_0}{a^4} \\
    \Rightarrow \rho_b = -\frac{4z^3P_0}{a^4} \\
\end{lefteqs}

\textbf{ב)}

\begin{lefteqs}
    Q_{total} = \iint_{z=0} \sigma_b dS + \iint_{z=a} \sigma_b dS + \iint_{x=0} \sigma_b dS + \iint_{x=a} \sigma_b dS + \iiint_V \rho_b dV \\
    = \int_0^a \int_0^a \frac{x}{2a}P_0 dxdy + \int_0^a \int_0^a(1-\frac{x}{2a})P_0 dxdy  + \int_0^a \int_0^a -P_0 dydz \int_0^a \int_0^a P_0 dydz \\
    + \int_0^a \int_0^a \int_0^a -\frac{4z^3P_0}{a^4} dxdydz \\
    = P_0 \frac{a}{4} + \frac{3P_0 a^2}{4} - P_0 a^2 + P_0 a^2 - a^2 P_0 = 0
\end{lefteqs}

זו התוצאה שציפינו לה, פולריזציה בחומר רק מפזרת את המטען ולא יוצרת מטען נוסף (לכל דיפול יש צד חיובי וצד שלילי)
ולכן סך המטענים המושרים תמיד יהיה 0.


\vspace{3em}
\hrule
\vspace{2em}

\textbf{\large שאלה 5}
(פרק 3 שאלה 5) 

נתון חומר בעל צפיפות $3.3 \cdot 10^{28}$ אטומים למטר בלישית.
פועל שדה מגנטי קבוע בגודל $H=4\frac{A}{m}$ על החומר המקבלים מכל אטום מומנט מגנטי ממוצע של $m=1.2 \cdot 10^{-24}$

\textbf{א)}

\begin{lefteqs}
    M=n \cdot m = 3.3 \cdot 10^{28} \cdot 1.2 \cdot 10^{-24} = 3.96 \cdot 10^4 \frac{A}{m} \\
    \chi_m = \frac{M}{H} = \frac{3.96 \cdot 10^4}{4} = 9.9 \cdot 10^3 \\
    \mu = \mu_0 (1 + \chi_m) = 4\pi \cdot 10^{-7}(1+9.9 \cdot 10^3) = 0.01244 \frac{H}{m}
\end{lefteqs}

\textbf{ב)}
\begin{lefteqs}
    B = \mu H = 0.01244 \cdot 4 = 0.4976 T
\end{lefteqs}


\vspace{3em}
\hrule
\vspace{2em}

\textbf{\large שאלה 6}
(פרק 3 שאלה 11)

באזור $y<0$ שורר ריק, וב $y>0$ יש מגנטיזציה שגודלה

$\vec{M} = M_0 \cos(\beta x) \exp{-\alpha y} \hat{z}$

נמצא את הפוטנציאל המגנטי בכל נקודה במרחב.

אין זרמין חיצוניים ($J_{ext}=0$) ולכן $\nabla^2 \phi_M = \vec{\nabla} \cdot \vec{M}$/

כאשר עבור $y<0$ מתקיים $M=0$ ועבור $y>0$ $M$ נתון בשאלה.

נקבל משוואת לפלאס ב $y<0$ ומשוואת פואסון ב $y>0$:

\begin{lefteqs}
    \nabla^2 \phi_M^- = 0 \\
    \nabla^2 \phi_M^+ = \vec{\nabla} \cdot \vec{M} = -\alpha M_0 \cos(\beta x) e^{-\alpha y}
\end{lefteqs}

עבור $y<0$ הפוטנציאל צריך לדעוך כאשר $y \to -\infty$ ולכן:
\begin{lefteqs}
    \phi_M^- = (A\sin(kx) + B\cos(kx)) \cdot (Ce^{ky} + De^{-ky}) \\
\end{lefteqs}

כאשר $D=0$ כדי לאפשר דעיכה של הפוטנציאל ב $y \to -\infty$:
\begin{lefteqs}
    \phi_M^- = (A\sin(kx) + B\cos(kx)) \cdot e^{ky}
\end{lefteqs}

עבור $y>0$ הפונציאל דועך כאשר $y \to \infty$, הפתרון יהיה הפתרון הכללי פלוס פיתרון פרטי:

\begin{lefteqs}
    \phi_M^+ = (C\sin(kx) + D\cos(kx)) \cdot e^ {-ky} + \phi_p
\end{lefteqs}

נחפש פתרון פרטי מהצורה: $\phi_p = E \cos(\beta x) e^{-\alpha y}$

\begin{lefteqs}
    \nabla^2 \phi_p = \vec{\nabla} \cdot \left[ E \beta \sin(\beta x)e^{-\alpha y}\hat{x} - E \alpha \cos(\beta x) e^{-\alpha y} \hat{y} \right] \\
    \nabla^2 = E \beta^2 \cos(\beta x)e^{-\alpha y} + E\alpha^2\cos(\beta x)e^{-\alpha y} = E\cos(\beta x)e^{-\alpha y}(\alpha^2 - \beta^2) \\
    E(\alpha^2 - \beta^2)\cos(\beta x)e^{-\alpha y} = -\alpha M_0 \cos(\beta x) \cdot e^{-\alpha y} \\
    \Rightarrow E(\alpha^2 - \beta^2) = -\alpha M_0 \Rightarrow E = \frac{-\alpha M_0}{\alpha^2 - \beta^2} \\
    \Rightarrow \phi_p = \frac{-\alpha M_0}{\alpha^2 - \beta^2} \cos(\beta x) e^{-\alpha y}
\end{lefteqs}

פונקצית הפוטנציאל צריכה להיות רציפה ב $y=0$, כלומר $\phi_M^-|_{y=0} = \phi_M^+|_{y=0}$:

\begin{lefteqs}
    A\sin(kx) + B\cos(kx) = C\cdot sin(kx) + D\cdot \cos(kx) -\frac{\alpha M_0}{\alpha^2 - \beta^2}\cdot \cos(\beta x) \\
    A = C \, k_1 = k_2 = \beta, B=D-\frac{\alpha M_0}{\alpha^2 - \beta^2} \\
\end{lefteqs}
שנית, $B$ הנורמלי צריך להיות רציף סביב $y=0$, מכיוון של $\vec{M}$ יש רק רכיב $y$:

\begin{lefteqs}
    B_y = mu_0 (H_y + M_y) \\
    \forall y<0 \, , M = 0 \, , B_y^- = \mu_0 H_y = mu_0 \frac{-\partial \phi_{M,y}^-}{\partial y} \\
    \forall y>0 \, , M = M_0 \cos(\beta x) e^{-\alpha y} \hat{z} \, , B_y^+ = \mu_0 (\frac{-\partial \phi_{M,y}^+}{\partial y} + M_0 \cos(\beta x)e^{-\alpha y}) \\
    B_y^-|_{y=0} = B_y^+|_{y=0} \Rightarrow \\
    -\beta(A\sin(\beta x) + B\cos(\beta x)) = \beta(C\sin(\beta x) + D\cos(\beta x)) - \frac{\alpha^2 M_0}{\alpha^2-\beta^2}\cos(\beta x) + M_0\cos(\beta x) \\
    -\beta A=\beta C \Rightarrow A=-C \Rightarrow A=C=0 \\
    \beta B = \beta D - \frac{\alpha^2 M_0}{\alpha^2 - \beta^2} + M_0 = \beta D + M_0 (1-\frac{\alpha^2}{\alpha^2 - \beta^2}) \\
    \beta D + M_0 (\frac{\alpha^2 - \beta^2 -\alpha^2}{\alpha^2 - \beta^2}) = \beta D - \frac{M_0 \beta^2}{\alpha^2 - \beta^2} \\
    B=D-\frac{\alpha M_0}{\alpha^2 - \beta^2}\\
    -\beta D _ \frac{\alpha \beta M_0}{\alpha^2 - \beta^2} = \beta D - \frac{M_0 \beta^2}{\alpha^2 - \beta^2} \\
    \Rightarrow 2\beta D = \frac{M_0 \beta (\alpha - \beta)}{\alpha^2 - \beta^2} = \frac{M_0 \beta(\alpha - \beta)}{(\alpha - \beta)(\alpha + \beta)} = \frac{M_0 \beta}{\alpha + \beta} \\
    \Rightarrow D = \frac{M_0}{2(\alpha + \beta)}
    B = \frac{M_0}{2(\alpha + \beta)} - \frac{\alpha M_0}{\alpha^2 -\beta^2} = \frac{M_0}{\alpha+\beta}(\frac{1}{2} - \frac{alpha}{\alpha-\beta}) = \frac{M_0}{\alpha + \beta}(\frac{\alpha-\beta-2\alpha}{2(\alpha-\beta)}) \\
    \Rightarrow B=\frac{M_0}{\alpha+\beta} \cdot (\frac{-(\alpha+\beta)}{2(\alpha-\beta)}) = \frac{-M_0}{2(\alpha-\beta)} \\
    \phi_M = \begin{cases}
    \frac{-M_0}{2(\alpha-\beta)} \cos(\beta x) e^{\beta y} & y<0 \\
    \frac{M_0}{2(\alpha+\beta)} \cos(\beta x) e^{-\beta y} - \frac{\alpha M_0}{\alpha^2 - \beta^2} \cos(\beta x) e^{-\alpha y} & y>0
    \end{cases}
\end{lefteqs}

\vspace{3em}
\hrule
\vspace{2em}

\textbf{\large שאלה 7}
(פרק 3 שאלה 16)

כדור שעליו צפיפות מטען משטחית $\sigma$ סובב סביב צירו (מקביל לציר $z$) בתדירות זוויתית $\omega$.

\textbf{א)}
נמצא את השדה המגנטי במרכז הכדור

ללא הגבלת הכלליות נניח כי הכדור ממוקם בראשית הצירים. נניח רדיוס $R$.

הזרם החשמלי המייצר את השדה המגנטי הוא כתוצאה מסיבוב הכדור, כלומר תנועת המטענים בתנועה הסביבות של הכדור.

נשים לב גם כי כל נקודה על הכדור נעה במהירות אחרת.

נסתכל על הכדור בתור הרבה טבעות, שכל אחת מסתובבת סביב ציר $z$, עם עובי אינפיניטסימלי.
רדיוס כל טבעת  הוא $R\sin(\theta)$ והמרחק של כל בעת ממרכז הכדור הוא $R\cos(\theta)$.

עבור טבעת בעובי $Rd\theta$:
\begin{lefteqs}
    dq = \sigma \cdot 2\pi R \sin(\theta) \cdot Rd\theta 
\end{lefteqs}
כלומר הזרם על כל טבעת הוא:
\begin{lefteqs}
    dI = \frac{dq}{T} = \frac{dq}{\frac{2\pi}{\omega}} = \omega \sigma R^2 \sin(\theta) d\theta
\end{lefteqs}
כידוע השדה המגנטי שיוצרת לולאת זרם הוא:
\begin{lefteqs}
    B = \frac{mu_0 I a^2}{2(a^2+d^2)^{\frac{3}{2}}}
\end{lefteqs}
נציב את הערכים:
\begin{lefteqs}
    \int_0^\pi \frac{mu_0 \omega \sigma R^2 \sin(\theta) \cdot R^2 \sin^2(\theta)}{2(R^2\sin^2(\theta) + R^2\cos^2(\theta))^{\frac{3}{2}}} d\theta \\
    = \omega \mu_0 \sigma \frac{R^4}{2R^3} \int_0^\pi \frac{\sin^3(\theta)}{sin^2(\theta) + \cos^2(\theta)} d\theta = \frac{\omega \mu_0 \sigma R}{2} \int_0^\pi \sin^3(\theta) d\theta \\
    = \frac{\omega \mu_0 \sigma R}{2} \left[-\frac{1}{3}\cos(\theta)(\sin^3(\theta) + 2) \right]_0^\pi \\
    = \frac{\omega \mu_0 \sigma R}{2} \left(\frac{1}{3}(0+2) + \frac{1}{3}(0+2) \right) = \frac{2\omega \mu_0 \sigma R}{3}
\end{lefteqs}

את הכיוון נקבע ע"פ כלל יד ימין:
\begin{lefteqs}
    \vec{B} = \frac{2\omega \mu_0 \sigma R}{3} \hat{z}
\end{lefteqs}

\textbf{ב)}
נמצא את השדה המגנטי בסדר המוביל רחוק מאוד ממרכז הכדור

כאשר הנקודת יחוס רחוקה מאוד ממקור הזרם, נעשה פיתוח מולטיפולי.

כאשר אין מונופולים מגנטיים ולכן הסדר הוא דיפול מגנטי.

עבור לולאת זרם, $\vec{m} = IS \hat{z}$  ולכן:
\begin{lefteqs}
    dm=dI \cdot A = \omega \sigma R^2 \sin(\theta) d\theta = \pi (R\sin(\theta))^2 \\
    m=\int_0^\pi dm = \int_0^\pi \omega \sigma \pi R^4\sin^3(\theta) d\theta = \omega \sigma \pi R^3 \frac{4}{3}
\end{lefteqs}
ע"פ 3.63

\begin{lefteqs}
    \vec{B} = \frac{mu_0 m}{4\pi r^3}(2\cos(\theta)\hat{r} + \sin(\theta)\hat{\theta}) \\
    \vec{B} = \frac{4\mu_0 \omega \sigma R^4}{3r^3} \cdot (2\cos(\theta)\hat{r} + \sin(\theta)\hat{\theta})
\end{lefteqs}

\end{document}